\documentclass[aspectratio=169]{ctexbeamer}
\usepackage[T1]{fontenc}
\usepackage{latexsym,amsmath,xcolor,multicol,booktabs,calligra}
\usepackage{graphicx,listings,stackengine}
\usepackage{hyperref}

\usepackage{PekingU}

% 去掉每节/子节前面的自动目录页
\AtBeginSection[]{}
\AtBeginSubsection[]{}

% ===== 每节课改这里 =====
\author{Cap1taL}
\title{一类 DP 问题选讲}
\date{2026年8月}
% ========================

% 代码高亮设置
\definecolor{deepblue}{rgb}{0,0,0.5}
\definecolor{deepred}{rgb}{0.6,0,0}
\definecolor{deepgreen}{rgb}{0,0.5,0}
\definecolor{halfgray}{gray}{0.55}

\lstset{
    basicstyle=\ttfamily\small,
    keywordstyle=\bfseries\color{deepblue},
    emphstyle=\ttfamily\color{deepred},
    stringstyle=\color{deepgreen},
    numbers=left,
    numberstyle=\small\color{halfgray},
    rulesepcolor=\color{red!20!green!20!blue!20},
    frame=shadowbox,
}

\begin{document}

\kaishu

% 封面
\begin{frame}
    \titlepage
\end{frame}

% 目录
% \begin{frame}{目录}
%     \tableofcontents[sectionstyle=show,subsectionstyle=show/shaded/hide,subsubsectionstyle=show/shaded/hide]
% \end{frame}

% ===== 从这里开始写你的内容 =====

\section{我 D 了好多好多 P}

\subsection{Codeforces 1796E Colored Subgraphs}

\begin{frame}{Codeforces 2001E2 Deterministic Heap (Hard Version)}
      给树做某种链剖分。你可以自由选择根和链剖分的策略，使得最短链最长
      
      $n\le 2\times 10^5$
\end{frame}

\begin{frame}{Solution}
    \begin{itemize}
        \item 先暴力枚举根节点，然后树形 DP。
        \item 考虑根节点怎么转移，显然我们应该延伸最短的那根树链，否则答案不优
        \item 因此我们设 $f_u$ 表示 $u$ 在上述策略下的可能最短链长
        \item 同时，又由于其他没有被我们延伸的链已经结束了，因此我们还需要知道他们的最小值，用于更新答案
        \item 因此我们还需要维护次小值 $g_u$，也就是 “没有机会继续被延伸” 的最小链
        \item 现在来处理换根的情况，注意到我们可能会出现递归进入次小值链的情况，所以我们还需要维护第三小值，用于次小值失效时作为新的次小值进行换根 DP。
        \item 复杂度 $O(n)$，如果用接受 $O(n\log n)$ 可以用 multiset 省去很多分类讨论。
    \end{itemize}
\end{frame}

\subsection{洛谷 P9746 合并序列}
\begin{frame}{洛谷 P9746 合并序列}
    给定序列 $a$。每次操作选择 $i<j<k$ 满足 $a_i\oplus a_j\oplus a_k=0$，将 $a_i..a_k$ 替换为它们的异或和。判断能否使序列长度为 $1$，并输出方案。
    
    $n\le 500,\ a_i<512$
\end{frame}

\begin{frame}{Solution}
    \begin{itemize}
        \item 注意到操作不会改变序列的异或和
        \item 区间 DP，直接枚举出三个子区间，复杂度 $O(n^6)$
        \item 发现每次都枚举中间那个区间非常蠢，设 $g_{l,r,k}$ 表示 $[l,r]$ 内是否存在可以被缩起来的异或和为 $k$ 的区间，其本身的更新的 $O(n^3)$ 的。优化到 $O(n^4)$
        \item 为了进一步优化 DP 的复杂度，我们再定义 $h_{l,r,k}$ 表示 $[l,r]$ 内是否存在 $[l,a],[b,c]$ 满足二者各自能缩起来，且异或和为 $k$
        \item 由此一来 DP 可以优化到 $O(n^3)$，但 $h$ 数组本身的求解是 $O(n^4)$ 的
    \end{itemize}
\end{frame}

\begin{frame}{Solution}
    \begin{itemize}
        \item 发现 $g,h$ 的定义是冗余的，以 $g$ 为例，若 $g_{l,r,k}$ 为真，那么 $g_{l,r+1,k}$ 也为真
        \item 更改定义，设 $g_{l,k}$ 表示最小的 $r$ 使得原来的 $g_{l,r,k}$ 成立，$q$ 同理
        \item 具体地，如何更新二者？若我们已有 $f_{l,r}$ 为真，那么首先需要枚举 $L\le l$ 来更新 $g_{L,k}$，再枚举一个异或和 $S$ 代表右侧区间，来更新 $q_{l,S}$
        \item 复杂度 $O(n^3)$，思路很正的一道区间 DP 优化
        \item 注意枚举顺序，应该采用倒序枚举左端点而不是区间长度递增的顺序。这是由我们的更新所决定的 
    \end{itemize}
\end{frame}

\subsection{洛谷 P7718 Equalization}
\begin{frame}{洛谷 P7718 Equalization}
    给定序列 $a[1..n]$。每次操作选 $l,r,x$，将 $a_l..a_r$ 均加上 $x$。求使所有元素相等的最少操作次数及方案数。
    
    $n\le 18$
\end{frame}

\begin{frame}{Solution}
    \begin{itemize}
        \item 好牛的状压
        \item 看到区间加，我们上来就是一个差分 $b_i$，操作变为在二者之间转移数的数量，或者做单点修改
        \item 那么显然就有了一个 $n-1$ 次的操作方案，每次都无脑转移给 $b_1$。对标原题大概是每次都使得一个数和 $a_1$ 一样
        \item 尝试减少操作次数。我们的操作目标实际上是使得除了 $b_1$ 之外都为 $ 0 $，那么一旦存在一个 $[2,n]$ 的子集，内部的 $b_i$ 本身为 $1$，我们就可以在子集内部操作，节省一次操作
        \item 问题转化为如何找尽量多的这种子集。设 $f_S$ 表示 $S$ 集合的和为 $ 0 $ 时最多划分成多少个和为 $ 0 $ 的子集
        \item 预处理一下可以做到 $O(3^n)$
    \end{itemize}
\end{frame}

\begin{frame}{Solution}
    \begin{itemize}
        \item 操作次数解决了，问题是如何计算方案数。
        \item 考虑图论建模，一次数的转移可以看做两点之间连边，一次单点修则是自环
        \item 那么不难发现我们的图形如若干个树，以及一个带自环的树
        \item 因为我们的边刻画的是数的转移，或者说数的流动，因此自然一个连边方案可以唯一确认所有边权
        \item 所以不需要考虑边权，问题转化为连边的方案。一个树的连边方案是 $|V|^{|V|-2}$，一个自环树的方案是 $2\times|V|\times|V|^{|V|-2}$
        \item 之后加加乘乘就可以算出答案，复杂度 $O(3^n)$
        \item 可以再做更精细的转化，搞到 $O(n^22^n)$。不过没必要。
    \end{itemize}
\end{frame}

\subsection{AtCoder agc049\_d Convex Sequence}
\begin{frame}{AtCoder agc049\_d Convex Sequence}
    长度为 $N$ 的非负整数序列 $A$，满足：
    \begin{itemize}
        \item $\displaystyle\sum_{i=1}^N A_i = M$
        \item 对任意 $2\le i\le N-1$，$2A_i \le A_{i-1}+A_{i+1}$（凸性）
    \end{itemize}
    求序列个数，对 $10^9+7$ 取模。
    
    $N,M\le 10^5$
\end{frame}

\begin{frame}{Solution}
    \begin{itemize}
        \item 要求我们计数一些凸数列的数量
        \item 自然地从最低点劈开，只关注一侧单调的情况
        \item 发现凸数列一定可以由一段后缀加 $[1,2,3,4,5..]$ 得到。证明考虑二阶差分即可
        \item 那么我们问题转化为选若干个后缀，执行若干次上述加操作，可以把区间操作抽象成物品，于是有了一个 $O(nm)$ 的背包
        \item 立刻发现这个后缀加的和很可能超过 $m$，只有 $O(\sqrt{m})$ 级别的区间是有用的，因此复杂度变成 $O(n\sqrt{m})$ 就对了
        \item 现在考虑变成两侧单调的情况，发现是对称的，我们在另一侧添加一个数，等价于左侧多了一种前缀区间操作，右侧少了一种后缀区间操作，类似于添加和删除物品
        \item 由于是计数背包，我们可以任意地删除物品
        \item 复杂度 $O(n\sqrt{m})$
    \end{itemize}
\end{frame}

\subsection{AtCoder abc261\_g Replace}
\begin{frame}{AtCoder abc261\_g Replace}
    给定字符串 $S,T$ 和 $K$ 种操作：将单个字符 $C_i$ 替换为字符串 $A_i$，代价 $1$。求将 $S$ 变为 $T$ 的最小总代价，或输出 $-1$。
    
    $|S|\le|T|\le 50,\ K\le 50$
\end{frame}

\begin{frame}{Solution}
    \begin{itemize}
        \item 观察到我们似乎在“膨胀” $S$ 串
        \item 这提示我们，每个 $S$ 里的字符，都对应了 $T$ 的一个子串
        \item 两个串匹配，考虑 $DP_{i,j}$ 表示我们在匹配 $S$ 的 $[1,i]$ 和 $T$ 的 $[1,j]$，这一部分是简单的
        \item 现在我们需要求出 $f_{l,r,c}$ 表示字符 $c$ 变成 $T$ 中 $[l,r]$ 的最小代价，来辅助 $dp_{i,j}$ 的转移。
    \end{itemize}
\end{frame}

\begin{frame}{Solution}
    求出 $f_{l,r,c}$
    \begin{itemize}
        \item 首先我们可能会有，从单字符变成单字符的操作，这一部分类似最短路，用弗洛伊德处理掉这类转移
        \item 然后考虑 $c$ 一定变成了一个 $A_i$。我们设 $g_{i,j,l,r}$ 表示 $A_i$ 的 $[1,j]$ 最少几步变成 $T$ 的 ${l,r}$
        \item 那么自然可以用 $g_{i,|A_i|,l,r}$ 来更新 $f_{l,r,c}$（要求 $c$ 可以一步变为 $A_i$）
        \item 至于 $g$ 自身的转移也是直接的，依旧类似双串匹配，有 $g_{i,j,l,r}=\min_{l\le k\le r}(g_{i,j-1,l,k}+f_{k+1,r,c})$
        \item 看起来好像转移还是有环，但是按照区间从短到长更新就好了
        \item 复杂度 $O(|T|^5)$ 可以通过，瓶颈在于 $g$ 的更新
    \end{itemize}
\end{frame}


\subsection{Codeforces 2001E2 Deterministic Heap (Hard Version)}

\begin{frame}{Codeforces 2001E2 Deterministic Heap (Hard Version)}
      大小为 $2^n-1$ 的完美二叉树（最大堆），初始全 $0$。每次操作选择一个节点 $v$，将根到 $v$ 路径上所有节点 $+1$。
      
      若某堆的前两次 $\mathrm{pop}$ 操作均确定（每次选择子节点时值不相等），称其为确定性最大堆。求恰好 $k$ 次操作后能得到多少种不同的确定性最大堆。
      
      $2\le n\le 100$，$1\le k\le 500$，$\sum n\le 100$，$\sum k\le 500$，$p$ 为质数
\end{frame}

\begin{frame}{Solution}
    \begin{itemize}
        \item 考虑两次 pop 操作的形态，发现第二次 pop 形如：先走一段第一次 pop 的操作，再拐弯去别的子树，此后相当于第一次 pop
        \item 这启示我们类似 "分段" 地 DP
        \item 考虑第二次 pop 时最为特殊的分叉点 $u$，我们不妨假设第一次 pop 去了它的左儿子，第二次却去了右儿子
        \item 这要求它的左儿子的两个儿子，均比它的右儿子小，这样才会拐弯
        \item 反之亦然，如果两次 pop 都去了左儿子，则要求左儿子的两个儿子，存在一个比它的右儿子大
    \end{itemize}
\end{frame}

\begin{frame}{Solution}
    状态设计
    \begin{itemize}
        \item 首先计算 $ 0 $ 次 pop 的情况，只有操作合法的要求
        \item 设 $f_{i,k}$ 表示子树高为 $i$，根节点值是 $k$ 的方案数
        $$
            f_{i,k}=\sum_{x+y\le k;x\neq y}f_{i-1,x}\times f_{i-1,y}
        $$
    \end{itemize}
\end{frame}

\begin{frame}{Solution}
    状态设计
    \begin{itemize}
        \item 再来计算一次 pop。
        \item 设 $g_{i,k,v}$ 表示树高为 $i$ 根节点为 $k$，两个儿子都不大于 $v$ 的方案数
        $$
            g_{i,k,v}=2\times \sum_{x+y\le k;x<y\le v}f_{i-1,x}\times g_{i-1,y,y}
        $$
        \item 其中的系数 $2$ 表示不知道哪一侧更大
    \end{itemize}
\end{frame}

\begin{frame}{Solution}
    状态设计
    \begin{itemize}
        \item 再来计算两次 pop。
        \item 设 $h_{i,k,v}$ 表示树高为 $i$ 根节点为 $k$，两个儿子的最大值不小于 $v$ 的方案数
        $$
            h_{i,k,v}=2\times \sum_{x+y\le k;x<y;v\le y}\left(  g_{i-1,y,x-1}\times g_{i-1,x,x}  + h_{i-1,y,x+1}f_{i-1,x}  \right)
        $$
        \item 其中的系数 $2$ 表示不知道哪一侧更大
        \item 第一项表示分叉向两侧，第二项表示同侧删两次。各自有各自的值大小要求，和之前的分析保持一致
    \end{itemize}
\end{frame}


% ===== 结束页 =====

\begin{frame}
    \begin{center}
        {\Huge\calligra Thanks!}
    \end{center}
\end{frame}

\end{document}
